\begin{exercise}[Median -- Projection in \(L^1\) \Coffeecup\Coffeecup]
	We define the median of \(X\) to be the set
    \[
        \mSym[X] \coloneq \set[\Big]{m \in \real : \prob{X \ge m}\ge \tfrac12, \prob{X\le m} \ge \tfrac12}.
    \]
	\begin{enumerate}[label=(\alph*), noitemsep]
        \item\label{it: (median) interval}
        Show that if \(a,b\in \mSym[X]\), then \((a,b) \subseteq \mSym[X]\).
        That is \(\mSym[X]\) is an interval!

        \item\label{it: (median) interval boundaries}
        The next goal is to show that the median interval is closed
        and characterize the upper and lower bound. Prove for \(F(x) = \prob{X\le x}\)
        \begin{alignat*}{2}
            \max \mSym[X] &= \inf \set{x \in \real : F(x) > \tfrac12}
            &&\eqcolon \overline{M},
            \\
            \min \mSym[X] &= \inf \set{x \in \real : F(x) \ge \tfrac12} \overset{\text{def.}}= F^{\leftarrow}(\tfrac12)
            &&\eqcolon \underline{M}.
        \end{alignat*}
        This has two components: \(\overline{M}\) and \(\underline{M}\) are the boundaries of \(\mSym[X]\)
        and they are contained in \(\mSym[X]\).

        \textbf{Hint:} Use \(\prob{X < x} = \prob{\bigcup_{k\in \nat}\set{X \le x - \tfrac1k}}\) and the continuity of measure.

        \item\label{it: characterization of distribution function}
        Prove that
        \[
            \begin{cases}
                F(x) > \tfrac12 & x > \overline{M} 
                \\
                F(x) = \tfrac12 & x \in (\underline{M}, \overline{M})
                \\
                F(x) < \tfrac12  & x < \underline{M}
            \end{cases}
        \]
        In particular, if the distribution function \(F\) jumps over
        \(\tfrac12\) such that there exists no \(x\) with \(F(x) = \frac12\),
        then \(\underline{M} = \overline{M}\). For the case
        \(\underline{M} \neq \overline{M}\) prove \(F(\underline{M})= \tfrac12\).

		\item\label{it: (median) projection in L1}
        Assuming \(X\in \cL^1\) prove that the median minimizes the \(L^1\) error, i.e.
        \[
            \mSym[X] =\arg\min_{m} \E{\abs{X-m}}.
        \]
		% If you wish to prove the general
		% case, note that \(f(m)=2F_X(m)-1\) is right-continuous and monotonously
		% increasing. Monotonicity also implies it can only have a countable number
		% of discontinuities.

 		\textbf{Hint:}
            Use \(\E{X}=\int_0^\infty \prob{X\ge t}dt = \int_0^\infty \prob{X>t}dt\)
            for \(X\ge 0\) to prove
            \[
                \E{\abs{X-m_2}} - \E{\abs{X-m_1}}
                = \int_{m_1}^{m_2} \underbrace{\prob{X\le x} - \prob{X>x}}_{=2F(x)-1}dx
            \]
	\end{enumerate}
\end{exercise}

\begin{solution}
\begin{enumerate}[wide=0pt]
    \item[\ref{it: (median) interval}]

    \item[\ref{it: (median) interval boundaries}]
    Since the distribution function \(F\) is right-continuous we have
    \[
        \prob{X \le \overline{M}} = F(\overline{M}) = \lim_{x \downarrow \overline{M}} F(x) \ge \tfrac12
    \]
    as all \(x \ge \overline{M}\) have the property \(F(x) > \tfrac12\) by
    definition of \(\overline{M}\). For the other direction we use the
    continuity of measure:
    \[\begin{aligned}
        \prob{X \ge \overline{M}}
        &= 1 - \prob{X < \overline{M}}
        \\
        &= 1- \prob[\Big]{\bigcup_{k\in \nat} \set{X \le \overline{M} - \tfrac1k}}
        \\
        &= 1 - \lim_{k\to \infty} F(\overline{M}-\tfrac1k)
        \\
        &\ge 1- \tfrac12 = \tfrac12
    \end{aligned}
    \]
    as any \(x< \overline{M}\) has the property \(F(x) \le \tfrac12\).
    Consequently, \(\overline{M} \in \mSym[X]\).
    As we have not used the strict inequalities due to the limits, the
    same arguments work for the minimum \(\underline{M} = \inf\set{x\in
    \real F(x) \ge \tfrac12}\). We only replace \(F(x) > \tfrac12\) by
    \(F(x) \ge \tfrac12\) for all \(x \ge \underline{M}\) for the first
    part and use \(F(x) < \tfrac12\) for all \(x < \underline{M}\) for
    the second. Now by definition of \(\overline{M}\) we have for any
    \(m > \overline{M}\) that \(F(m) > \tfrac12\) and therefore
    \[
        \prob{X \ge m} = 1 - F(m) < \tfrac12 \qquad \forall m > \overline{M}.
    \]
    This implies \(\overline{M} = \max \mSym[X]\).
    Similarly, we have for any \(m < \underline{M}\) that \(F(m) < \tfrac12\)
    and therefore \(\underline{M} = \min \mSym[X]\).

    \item[\ref{it: characterization of distribution function}]
    \(F(x) > \tfrac12\) for \(x>\overline{M}\) follows from the definition
    of \(\overline{M}\). Similarly \(F(x) < \tfrac12\) for \(x<\underline{M}\)
    follows from the definition of \(\underline{M}\). Now for any
    \(x \in (\underline{M}, \overline{M})\) we have \(x> \underline{M}\) and
    consequently \(F(x) \ge \tfrac12\). And since \(F(x) > \tfrac12\) would imply
    \(x > \overline{M}\) by definition of \(\overline{M}\) we must conclude
    \(F(x) = \tfrac12\). Finally if \(\underline{M} \neq \overline{M}\)
    then \(F(\underline{M}) = \tfrac12\) follows from right-continuity of \(F\).

    \item[\ref{it: (median) projection in L1}]
    We have
    \begin{align*}
        \E{\abs{X-m}_+}
        &= \int_0^\infty \underbrace{\prob{\abs{X-m}_+ > t}}_{
            \begin{aligned}
                &=\prob{X-m>t}\\
                &=\prob{X> t-m}
            \end{aligned}
        } dt\\
        &= \int_m^\infty \prob{X>s}ds
    \end{align*}
    And similarly
    \begin{align*}
        \E{(X-m)_-}
        &= \int_0^\infty \underbrace{\prob{(X-m)_- \ge t}}_{
            \begin{aligned}
                &=\prob{m-X\ge t} \quad \mathrlap{(\forall t>0)}\\
                &=\prob{m-t\ge X}
            \end{aligned}
        } dt\\
        &= \int_{-\infty}^m \prob{X\le s}ds
    \end{align*}
    Put together we have
    \[
        \E{\abs{X-m}}
        = \int_{-\infty}^m \prob{X\le s}ds + \int_m^\infty \prob{X>s}ds.
    \]
    Consequently, we have for \(m_1 < m_2\)
    \[\begin{aligned}
        \E{\abs{X-m_2}} - \E{\abs{X-m_1}}
        &= \int_{m_1}^{m_2} \prob{X\le s}ds - \int_{m_1}^{m_2} \prob{X>s}ds
        \\
        &= \int_{m_1}^{m_2} \underbrace{\prob{X\le s} - (1-\prob{X\le s})}_{=2F(x) -1}ds
    \end{aligned}
    \]
    Using \ref{it: characterization of distribution function} we deduce
    \begin{alignat*}{3}
        \E{\abs{X-m}} - \E{\abs{x-\overline{M}}} &> 0
        & m &> \overline{M} 
        \\
        \E{\abs{X-m_2}} - \E{\abs{x-m_1}} &= 0
        \qquad & m_1, m_2 &\in (\underline{M}, \overline{M})
        \\
        \E{\abs{x-\underline{M}}} - \E{\abs{X-m}} &< 0
        & m &< \underline{M}
    \end{alignat*}
    Reordering implies the claim. \qedhere
\end{enumerate}
\end{solution}
